% Discussion for item 5.
% scripts/make_latex.py will NEVER overwrite this file -- edit it freely.

\paragraph{5.1 --- the sender reads back its own message.}
The outcome depends entirely on whether the sender still holds the read end, so
both cases were tested.

With the read end closed, as item 4 requires, the child's \texttt{read()} failed
at once with \texttt{EBADF} --- ``bad file descriptor''. The error is about the
\emph{handle}, not the pipe: \texttt{fd[0]} is no longer a valid descriptor in
that process, so the call never reaches the pipe to discover whether anything is
in it. Nothing was consumed, and the parent received the message normally.

With the read end left open, the child read its own 28-byte message back
successfully, and the parent then got \texttt{0} --- end of file --- with nothing
to show for it. A pipe is not addressed to anybody. It is one shared buffer, and
whichever process reads first takes the bytes. The convention that a pipe runs
one way is maintained purely by both sides closing the end they do not use; the
kernel does not enforce it. Note the parent's \texttt{read()} returned 0 rather
than blocking: by then the child had exited, so no write end remained open
anywhere, and the kernel reports end-of-file rather than making the reader wait
for data that can never arrive.

\paragraph{5.2 --- the receiver reads before the sender sends.}
The parent called \texttt{read()} at $t=0$ on an empty pipe and the call returned
at $t=3.004$~s, immediately after the child wrote. It did not fail and it did not
return 0; it blocked, and the process was simply set aside until data arrived.

This is the distinction the previous case turned on. An empty pipe that still has
a writer means ``no data \emph{yet}'', and the reader waits. An empty pipe with no
writers left means ``no data \emph{ever}'', and \texttt{read()} returns 0. The
same empty buffer produces opposite behaviour depending on whether any write end
is still open --- which is why closing the unused ends correctly matters so much.

\paragraph{5.3 --- several messages before a single read.}
The child made three separate \texttt{write()} calls of 32 bytes each. The parent
made one \texttt{read()} and received 96 bytes: all three messages, concatenated,
with nothing marking where one ended and the next began.

A pipe is a byte stream, not a queue of messages. The three writes did not create
three records; they appended bytes to one buffer. This is the essential difference
between a pipe and a message queue, which preserves message boundaries. If the
receiver needs its messages back as separate items, the two sides must agree on a
framing convention of their own --- a delimiter such as a newline, or a fixed-size
length prefix before each message. Here the messages happen to end in newlines, so
the receiver could split on those, but that only works because both sides
implicitly agreed on it.

\paragraph{In summary.}
A correct pipe user closes the end it does not need, so that end-of-file is
reported when the writers are really gone; expects \texttt{read()} to block rather
than to fail when data has not arrived yet; and does not assume that one
\texttt{write()} corresponds to one \texttt{read()}.
